\documentclass[12pt]{article}
\usepackage[brazil]{babel}
\usepackage[utf8]{inputenc}
\usepackage[T1]{fontenc}
\usepackage{amsmath, amssymb}
\usepackage{geometry}
\usepackage{graphicx}
\usepackage{xcolor}
\usepackage{enumitem}

\definecolor{maincolor}{RGB}{0,100,50}
\usepackage{fancyhdr}
\definecolor{forestgreen}{RGB}{0,100,50}
\usepackage[most]{tcolorbox}

\newtcolorbox{block}[1]{%
  colback=green!5,
  colframe=maincolor,
  coltitle=white,
  title=#1,
  fonttitle=\bfseries,
  boxrule=1.5pt,
  arc=3mm,
  left=6pt, 
  right=6pt,
  top=6pt,
  bottom=6pt
}

\geometry{margin=2.5cm}
\pagestyle{fancy}
\fancyhf{}

\renewcommand{\headrulewidth}{2pt}
\renewcommand{\footrulewidth}{2pt}

\renewcommand{\headrule}{\hbox to\headwidth{%
  \color{forestgreen}\leaders\hrule height \headrulewidth\hfill}}

\renewcommand{\footrule}{\hbox to\headwidth{%
  \color{forestgreen}\leaders\hrule height \footrulewidth\hfill}}

\setlength{\headheight}{32pt}

\fancyhead[L]{%
  % Certifique-se de que o arquivo de imagem esteja na mesma pasta
  \includegraphics[height=1.2cm]{Universidade_de_Rio_Verde.jpg}
}

\fancyhead[R]{%
  \raggedleft
  \textcolor{forestgreen}{%
    \small
    \textbf{Lista de Exercícios 3.}\\
    Me.~Sergio Souza Novak\\
    \textit{Lógica para Engenharia de Software}
  }
}

\fancyfoot[C]{%
  \textcolor{forestgreen}{\thepage}
}

\begin{document}

\vspace*{0.2cm}
\section*{Exercícios — Relações Binárias, Propriedades e Equivalência}

\begin{enumerate}

\item Sejam os conjuntos \(A = \{1, 2, 3\}\) e \(B = \{a, b\}\). Determine:
\begin{enumerate}[label=\alph*)]
    \item O produto cartesiano \(A \times B\).
    \item O produto cartesiano \(B \times A\).
    \item Os dois produtos são iguais? Justifique.
\end{enumerate}

\item Seja \(A = \{1, 2, 3, 4, 6\}\) e \(R\) uma relação em \(A\) definida por ``\(x\) divide \(y\)'', ou seja, \(x\,R\,y \Leftrightarrow x \mid y\).
\begin{enumerate}[label=\alph*)]
    \item Escreva os pares de \(R\) por extensão (liste todos os pares).
    \item Escreva a matriz relacional \(M_R\) associada a esta relação.
\end{enumerate}

\item Considere o conjunto \(A = \{1, 2, 3, 4\}\). Para cada uma das relações em \(A\) abaixo, verifique se são reflexivas, irreflexivas, simétricas, antissimétricas, assimétricas ou transitivas. Justifique sua resposta caso alguma propriedade falhe.
\begin{enumerate}[label=\alph*)]
    \item \(R_1 = \{(1,1), (2,2), (3,3), (4,4), (1,2), (2,1)\}\)
    \item \(R_2 = \{(1,2), (2,3), (1,3)\}\)
    \item \(R_3 = \{(1,1), (1,2), (2,1), (2,2), (3,4), (4,3)\}\)
    \item \(R_4 = \varnothing\) (a relação vazia)
\end{enumerate}

\item Defina e explique formalmente o que é uma \textbf{Relação de Equivalência}. Quais são as três propriedades que devem ser satisfeitas simultaneamente?

\item Determine se as relações abaixo, definidas sobre o conjunto dos números inteiros \(\mathbb{Z}\), são de equivalência:
\begin{enumerate}[label=\alph*)]
    \item \(a\,R\,b \Leftrightarrow a = b\).
    \item \(a\,R\,b \Leftrightarrow a \le b\).
    \item \(a\,R\,b \Leftrightarrow a\) e \(b\) têm a mesma paridade.
\end{enumerate}

\item Seja a relação de congruência módulo 4 em \(\mathbb{Z}\), definida por \(a \equiv b \pmod{4}\) se, e somente se, \(4 \mid (a - b)\). Sabendo que esta é uma relação de equivalência:
\begin{enumerate}[label=\alph*)]
    \item Determine quais são as classes de equivalência desta relação.
    \item Mostre que o conjunto dessas classes forma uma partição de \(\mathbb{Z}\).
\end{enumerate}

\item Defina \textbf{Partição} de um conjunto. Suponha que um conjunto \(A = \{a, b, c, d, e\}\) possui a partição \(\mathcal{P} = \{\{a, b\}, \{c\}, \{d, e\}\}\).
\begin{enumerate}[label=\alph*)]
    \item Verifique se a coleção cumpre os três requisitos formais de uma partição.
    \item Defina, descrevendo explicitamente os pares ordenados, a relação de equivalência \(R_\mathcal{P}\) induzida por esta partição.
\end{enumerate}

\item (Aplicação em Engenharia de Software). Em testes de sistema, aplica-se a técnica de \textit{particionamento de equivalência} das entradas. Suponha que um software recebe um valor de idade \(x \in \mathbb{N}\). Define-se a relação:
\(x \sim y \Leftrightarrow f(x) \text{ e } f(y) \text{ exibem o mesmo comportamento}\).
Considere que o sistema trata as idades diferentemente nestes intervalos: Criança (0 a 12), Adolescente (13 a 17) e Adulto (18 em diante).
\begin{enumerate}[label=\alph*)]
    \item Quantas classes de equivalência existem neste sistema?
    \item Como essa abordagem reduz o número de casos de teste que um engenheiro de software precisa escrever?
\end{enumerate}

\end{enumerate}

\end{document}
